\begin{figure*}[!t]
  \centering
  \begin{minipage}{\textwidth}
    \centering
    \textbf{(a)} Target-only analysis (numerosities 1-6)\par\smallskip
    \includegraphics[width=\textwidth]{../results/analysis-lda/static-change6/static_change6_all/figures-pairwise/temporal_model_rdms_all_compact_row-static_change6_all-pairwise.png}
  \end{minipage}
  \vspace{4pt}
  \begin{minipage}{\textwidth}
    \centering
    \textbf{(b)} Target-only no-1 analysis (numerosities 2-6)\par\smallskip
    \includegraphics[width=\textwidth]{../results/analysis-lda/static-change5/static_change5_all/figures-pairwise/temporal_model_rdms_all_compact_row-static_change5_all-pairwise.png}
  \end{minipage}
  \caption{Brain and model RDMs (0-700 ms). Brain RDMs show mean
  balanced accuracy (\%) across subjects. Model RDMs show predicted
  dissimilarity for each theoretical model. The RT model shows absolute
  differences in grand-mean reaction time. The Pixel model shows
  absolute differences in white-pixel area. In this stimulus set, total
  area is similar within numerosities 1-3 and within 4-6 but drops
  sharply between 3 and 4, so the Pixel model's largest discontinuity
  falls at 3-4. (a)~Target numerosities 1-6. Pairs involving
  numerosity~1 show the highest discriminability, giving the brain RDM
  an ANS-like gradient dominated by distance from~1
  (Table~\ref{tab:static-model-comparison}). In the boundary region,
  pairs 4-5 and 4-6 show slightly higher discriminability than pairs 2-4
  and 3-4, consistent with a boundary at 4-5 rather than 3-4.
  (b)~Target numerosities 2-6 (numerosity~1 excluded from both prime and
  target during training and evaluation). With numerosity~1 removed, the
  boundary structure becomes clearer. Pairs crossing the 4-5 boundary
  (4-5, 4-6) show higher discriminability than pairs within the PI range
  (2-4, 3-4), matching the Target PI(2-4) prediction
  (Table~\ref{tab:static-model-comparison}). The Pixel model predicts
  the opposite pattern, with highest dissimilarity across the 3-4
  boundary, and does not reach significance
  (Table~\ref{tab:static-model-comparison}).}
  \label{fig:static-model-rdms}
\end{figure*}
